\section{Maximum closure and its linear program}\label{sec:closure}

A \emph{closure} of a directed graph is a set of vertices containing all successors
of its members; with the arc convention of this project, \(i \to j\) meaning
``\(i\) requires \(j\)'', a closure is a set closed under prerequisites. The
maximum-weight closure problem is solved by a minimum cut in a network with
infinite-capacity arcs on the precedence relation, a reduction due to
\citet{Picard1976} and, in the selection form, to \citet{Rhys1970} and
\citet{Balinski1970}; \citet{JohnsonNiemi1983} and \citet{Gale1957} belong to
the same early line. The polyhedral fact that makes the reduction work --- that
the constraint matrix of the closure inequalities is an interval-free network
matrix, so the natural LP relaxation is integral --- is what allows this project
to work with the LP relaxation and still speak about closures.

\citet{PicardQueyranne1980} characterise the structure of \emph{all} minimum
cuts of a network, and \citet{PicardQueyranne1982} survey the applications of
minimum cuts to combinatorial problems, including closure. Both are used here
for a specific technical point rather than for background: the distinction
between the minimal and the maximal source side of a minimum cut, which is what
the canonical convention on tied ratios needs, is theirs.

On the algorithmic side, \citet{Hochbaum2001} revisits minimum cut and maximum
flow specifically in closure graphs, \citet{HochbaumChen2000} compare
implementations on the open-pit instance of the problem, and
\citet{Hochbaum2008} gives the pseudoflow algorithm, which is the method a
serious flow-based competitor would use today. \citet{UnderwoodTolwinski1998}
is worth singling out: they solve the ultimate pit problem with a network flow
algorithm derived from the \emph{dual}, and relate its steps to those of
Lerchs--Grossmann, so the idea of attacking maximum closure through dual
information is itself old. \citet{AmankwahEtAl2014} extend Picard's reduction
to the multi-period case on a time-expanded graph, and \citet{MeagherEtAl2014}
review the design of pushbacks, which is where the sequence of nested closures
produced by the canonical decomposition is actually consumed. \citet{AhujaEtAl1993} and
\citet{Schrijver2003} are the standard references for the surrounding network
flow and combinatorial optimisation material.

A generalisation worth recording, because it places the closure problem inside a
larger family rather than beside it: \citet{HochbaumQueyranne2003} minimise a
convex cost over a closed set, which specialises to maximum closure on one side
and to isotonic regression on the other. The same family is visited from the
other end by the maximum density subgraph, which is the ratio objective of
\cref{sec:ratio} without any precedence structure at all --- solved by flow in
\citet{Goldberg1984}, given its linear program by \citet{Charikar2000}, and
treated as one parametric cut by \citet{Hochbaum2024}.

Not everything in this line is a flow algorithm, and the exception matters
enough to have its own place: the algorithm of \citet{LerchsGrossmann1965}
maintains a \emph{forest} and pivots on it. It is discussed in \cref{sec:lg},
next to the recent forest-structured work, because that is where it does its
damage.

\paragraph{What this costs the claim.} Everything about the closure model, its
LP relaxation, its integrality and its solution by minimum cut is classical.
Nothing in \FS{} may be presented as new at this level, and the project's own
documents say so: the mathematical contract states the model and attributes it,
and the tutorial recalls the reduction only to fix notation.
